\section{Method}

\subsection{Overview}

This project defines a modular decoding workflow from multichannel sEMG and gesture annotations to five-finger prediction logits. The pipeline contains four stages: data harmonization, signal processing, model inference, and evaluation or deployment analysis. Figure~\ref{fig:pipeline} summarizes the workflow used throughout the study.

\begin{figure}[t]
  \centering
  \includegraphics[width=\linewidth,trim=35 120 35 70,clip]{pipeline_overview.png}
  \caption{Overview of the sEMG finger-intent decoding pipeline. Raw sEMG and gesture annotations are aligned, segmented, featurized, mapped into shared tensors, evaluated across multiple model families, and prepared for downstream deployment analysis.}
  \label{fig:pipeline}
\end{figure}

For notation, let a pre-segmented recording be $X \in \mathbb{R}^{T \times C}$, where $T$ is the number of raw time samples and $C=64$ is the channel count under the \physiomio\ path. After preprocessing and windowing, the project extracts a feature tensor $\mathcal{X} \in \mathbb{R}^{C \times W \times F}$, where $W$ is the number of windows and $F=12$ is the number of handcrafted descriptors per channel-window pair. Each sample is paired with a multilabel target vector $y \in \{0,1\}^5$ indicating intended activation of thumb, index, middle, ring, and little finger.

\subsection{Data harmonization and ingestion}

The implementation builds processed datasets directly from raw \physiomio\ parquet files. The loader searches the raw data tree, groups files by patient, aligns contiguous stretches of the same \texttt{movement\_type}, maps each label to a five-bit finger target, discards segments shorter than 200 samples, and applies the preprocessing pipeline to each usable segment. This converts gesture-level clinical recordings into segment-level multilabel examples.

The default configuration uses \physiomio\ recordings sampled at 2000 Hz and selects the impaired arm. Patients are partitioned into deterministic 70/10/20 train/validation/test splits, so samples from the same patient do not appear in multiple partitions. Processed samples are padded to a dataset-wide maximum window count for batched tensor operations.

\subsection{Signal preprocessing}

The preprocessing module is explicitly parameterized by a frozen configuration object in the implementation. Raw sEMG first passes through a fourth-order Butterworth band-pass filter with 20--450 Hz cutoffs and zero-phase forward-backward application. This stage suppresses low-frequency motion artifacts and high-frequency electrical noise while preserving time alignment across channels.

The filtered signal is then denoised with wavelet soft-thresholding. The default configuration uses a Symlet-4 basis, four decomposition levels, universal threshold selection, and a median-based noise estimate. The same preprocessing configuration is used for all model families, making downstream differences attributable to the representation and predictor rather than to separate signal-conditioning choices.

After denoising, each continuous segment is divided into 200 ms windows with 50\% overlap and no zero-padding. These settings balance temporal responsiveness against the sample count required for stable summary statistics.

\subsection{Window-level feature representation}

Each channel-window pair is summarized by twelve handcrafted descriptors spanning time and frequency domains. The time-domain set includes root mean square, mean absolute value, integrated EMG, waveform length, variance, zero crossings, slope sign changes, and Willison amplitude. The frequency-domain set includes mean frequency, median frequency, spectral entropy, and total power, with spectral quantities estimated from Welch periodograms using \texttt{nperseg = min(256, len(x))}~\citep{welch1967fft}. Degenerate spectra are zeroed explicitly to avoid unstable feature values. The resulting tensor has shape $(C, W, F)$, where $C$ is the number of channels, $W$ is the number of windows, and $F=12$ is the feature count.

\subsection{Tensor formatting and shared adapter}

The project standardizes cross-model comparison with one shared adapter. It transforms $(C, W, F)$ or $(N, C, W, F)$ tensors into $(N, W, C \times F)$ sequences, where windows act as time steps and all channel features are concatenated into one vector per step. For a single sample, the adapted sequence can be written as
\[
S = \left[s_1, \ldots, s_W\right]^\top \in \mathbb{R}^{W \times d},
\qquad
s_w = \mathrm{vec}\!\left(\mathcal{X}_{:,w,:}\right),
\qquad
d = C F = 768.
\]
This adapter creates a common input space for recurrent, convolutional, and graph-based predictors.

For the CNN branch, the adapted tensor is subsequently permuted to a channels-first layout. The CNN stack treats the $64 \times 12$ channel-feature grid as 768 input channels for temporal \texttt{Conv1d} layers, preserving the same upstream representation while allowing convolutions over window order.

\subsection{Predictor families}

The LSTM branch applies two stacked recurrent layers with hidden sizes 128 and 256, followed by a 128-unit fully connected layer and dropout before five output logits. Given adapted sequence input $s_t$, the recurrent dynamics follow the standard LSTM update
\[
\begin{aligned}
i_t &= \sigma(W_i s_t + U_i h_{t-1} + b_i), \\
f_t &= \sigma(W_f s_t + U_f h_{t-1} + b_f), \\
o_t &= \sigma(W_o s_t + U_o h_{t-1} + b_o), \\
\tilde{c}_t &= \tanh(W_c s_t + U_c h_{t-1} + b_c), \\
c_t &= f_t \odot c_{t-1} + i_t \odot \tilde{c}_t, \\
h_t &= o_t \odot \tanh(c_t),
\end{aligned}
\]
and the project uses the final hidden state from the second recurrent layer for classification. This branch emphasizes temporal accumulation across windows and serves as the main sequential baseline for the direct benchmark.

The GNN branch interprets each time window as a node in a graph. The implementation supports GCN, GraphSAGE, and GAT operators~\citep{kipf2017gcn,hamilton2017graphsage,velickovic2018gat}, and the benchmarked path constructs a dense graph over windows within each sample before graph-level pooling and readout. At a generic level, the node update can be written as
\[
h_i^{(\ell+1)} =
\psi_\ell\!\left(
h_i^{(\ell)},
\operatorname*{AGG}_{j \in \mathcal{N}(i)}
\phi_\ell\!\left(h_i^{(\ell)}, h_j^{(\ell)}\right)
\right),
\]
followed by graph pooling
\[
g = \operatorname*{POOL}_{i=1}^{W} h_i^{(L)}
\]
and a multilabel readout head. This branch injects an explicit structural prior: windows are related observations rather than only a flat sequence.

The CNN branch includes five student variants ranging from Nano to XLarge, each built around a $1 \times 1$ projection, temporal convolution blocks, adaptive pooling, and a compact fully connected head. If the channels-first representation is denoted by $\hat{S} \in \mathbb{R}^{d \times W}$, the student stack can be summarized as
\[
H^{(0)} = \rho\!\left(\operatorname{BN}(W_p * \hat{S} + b_p)\right),
\qquad
H^{(\ell+1)} = \rho\!\left(\operatorname{BN}(W_\ell * H^{(\ell)} + b_\ell)\right),
\]
where $*$ denotes 1D convolution over window order and $\rho$ is a ReLU nonlinearity. Adaptive pooling compresses the temporal axis before a multilayer perceptron produces the five logits. The student family spans approximately 54K parameters and 53 KB INT8 for Nano to 1.62M parameters and 1.58 MB INT8 for XLarge. The teacher family uses 1D ResNet-50, ResNet-101, and ResNet-152 variants with bottleneck residual blocks.

\subsection{Training, tuning, and evaluation protocol}

The training utilities unify optimization and reporting across model families. The current implementation uses binary cross-entropy with logits for five-label prediction, Adam-based optimization, plateau-triggered learning-rate reduction, early stopping, checkpointing, and saved metric curves. For logits $z \in \mathbb{R}^5$, the base multilabel objective is
\[
\mathcal{L}_{\mathrm{BCE}}
=
-\frac{1}{5}
\sum_{k=1}^{5}
\left[
y_k \log \sigma(z_k)
+ (1-y_k) \log \left(1-\sigma(z_k)\right)
\right].
\]
Evaluation reports exact-match accuracy, finger-level averaged accuracy, macro precision, macro recall, macro F1, macro AUROC, and macro AUPRC. The metrics module also stores per-finger scores and task-level curve data, which enables both aggregate and finger-wise analysis after training.

For the CNN students, Optuna searches architecture and training choices for each model size, and the optimization score combines validation accuracy with a latency-based penalty relative to a baseline timing reference~\citep{akiba2019optuna}. Trial-level early stopping is used during search. The system also supports teacher-ensemble distillation, per-finger threshold tuning, and latency benchmarking. For multilabel distillation, the implemented loss blends the hard target term with a temperature-scaled soft target from the teacher,
\[
\mathcal{L}_{\mathrm{KD}}
=
\alpha \, \mathcal{L}_{\mathrm{BCE}}(z_s, y)
+
(1-\alpha) T^2
\operatorname{BCELogits}\!\left(\frac{z_s}{T}, \sigma\!\left(\frac{z_t}{T}\right)\right),
\]
where $z_s$ and $z_t$ are student and teacher logits, $T$ is the distillation temperature, and $\alpha$ controls the hard-soft balance. A separate latency benchmark reports mean, median, and 95th-percentile inference time on a chosen device.

\subsection{Transfer-learning path}

In addition to direct impaired-arm benchmarks, the study evaluates two-stage CNN and LSTM training. Models are first pretrained on healthy-arm recordings and then finetuned on impaired-arm recordings while preserving the same five-finger multilabel objective. This tests whether healthy-limb structure can provide a useful initialization for paretic-limb decoding.
